\subsection{Witten--Reshetikhin--Turaev invariants} %\label{subsec:WRT}

Reshetikhin--Turaev~\cite{RT} constructed the invariant $\tau_{k}$ for a pair $(M, L)$ of a 3-manifold $M$ and a~framed link $L$ in $M$.
This invariant is parametrised by a positive integer $k$.
In our situation \mbox{$L = \varnothing$}, so $\tau_{k}$ is an invariant for 3-manifolds.
This invariant is called the $\SU(2)$ Witten--Reshetikhin--Turaev (WRT) invariant.
Turaev gives a detailed description~\cite{R}.

For plumbed manifolds, Gukov--Pei--Putrov--Vafa caluculated the WRT invariant expli\-ci\-tly~\cite[formula~(A.10)]{GPPV}:
\begin{align*}
	\tau_{k}(M_{3}(\Gamma))
	={}
	&\frac{\bm{e}(-\sigma/8) \zeta_{k}^{3\sigma/4}}{2(2k)^{V/2} \big(\zeta_{2k}-\zeta_{2k}^{-1}\big)}
	\sum_{n \in (\Z^{V} / 2k\Z^{V}) \smallsetminus (k\Z^{V} / 2k\Z^{V})}
	\prod_{v=1}^{V}\zeta_{k}^{w_{v}(n_{v}^{2} - 1) / 4}
	\\
	&\times\bigg( \frac{1}{\zeta_{k}^{n_{v}/2} - \zeta_{k}^{-n_{v}/2}} \bigg)^{\deg(v)-2}
	\prod_{(v, v') \in \mathrm{Edges}} \frac{\zeta_{k}^{n_{v}n_{v'}/2} - \zeta_{k}^{-n_{v}n_{v'}/2}}{2},
\end{align*}
where $ b_+ $ and $ b_- $ are the number of positive and negative eigenvalues counted with multiplicities of the linking matrix of a plumbed graph $\Gamma$,
$\sigma := b_+ - b_- $ is the signature,
$V$ is the number of vertices of $\Gamma$,
$ \deg(v) $ is the degree for an $v$th vertex,
and $w_{v}$ is a weight for an $v$th vertex.
We~use this formula in Section \ref{sec:WRT}.

\subsection{Homological blocks} %\label{subsec:homological_block}

In this section, we show a definition of homological blocks and express them as false theta functions.
Though any mathematical definitions of these invariants are not given for general cases, they provided a rigorous definition for a particular case in which 3-manifolds are plumbed.

A plumbed graph is a tree whose vertices are weighted by integers.
A 3-manifold is said to be plumbed if it is obtained by the surgery along a plumbed graph.
\begin{dfn}[{\cite[Section 3.4]{GPPV}}]
	Let $G$ be a plumbed graph with vertices $1, \dots, N$ and $\delta$ be as in Section \ref{sec:intro}.
	For the plumbed $ 3 $-manifold $M_{3}(G)$ and $ b \in ( H_1(M_3(G), \Z) + \delta)/ \{ \pm1 \} \cong \big(2\Z^{N} / 2W\big(\Z^{N}\big) + \delta\big) / \{ \pm1 \} $, the homological block $\widehat{Z}_{G, b} (q)$ is defined as follows:
	\begin{gather*}
	\widehat{Z}_{G, b} (q)
	=
	q^{-\sum_{v = 1}^{N} (w_{v} + 3)/4}
	\mathrm{v.p.} \int_{\abs{z_{v}}=1, v =1, \dots, N} \prod_{v = 1}^{N} \frac{{\rm d}z_{v}}{2\pi {\rm i} z_{v}} (z_{v} - 1/z_{v})^{2 - \deg(v)}
	\Theta_{-W, b}(z),
	\end{gather*}
	where $w_{v}$ is the weight of the vertex $v$, $ \mathrm{v.p.} $ denotes the Cauchy principal value, $\deg(v)$ is the degree of $v$, $W$ is the linking matrix of $\calL(G)$, and
	\begin{gather*}
	\Theta_{-W, b} (q; z_1, \dots, z_N) :=
	\sum_{\bm{l} = \trans{(l_1, \dots, l_N)} \in 2W (\Z^N) + b} q^{-\trans{\bm{l}} W^{-1} \bm{l}/4} \prod_{1 \le v \le N} z_v^{l_v}.
	\end{gather*}
\end{dfn}

If $ G $ is an H-graph $ \Gamma $, the above definition of $ \widehat{Z}_{G, b} (q) $ coincides with (\ref{eq:homological_block_for_H}).
\begin{dfn} %\label{dfn:false_theta}
	For $ \tau \in \bbH$, we define false theta functions $F_{Q_{+}, \veps}(\tau)$ and $F_{Q_{-}, \veps}(\tau)$ by
	\begin{gather*}
	F_{Q_{+}, \veps}(\tau) := \sum_{0 \le \gamma \in (2S)^{-1} (\Z^2)} \veps(\gamma) q^{Q_+(\gamma)}, \qquad
	F_{Q_{-}, \veps}(\tau) := \sum_{0 \le \gamma \in (2S)^{-1} (\Z^2)} \veps(\gamma) q^{Q_-(\gamma)},
	\end{gather*}
	where $ Q_+(m, n) := Q(m, n)$, $Q_-(m, n) := Q(m, -n) $.
\end{dfn}

In order to calculate radial limits of the homological block $\widehat{Z}_\Gamma (q)$, we need an expression of~$\widehat{Z}_\Gamma (q)$ as a sum of false theta functions given by Bringmann--Mahlburg--Millas.

\begin{prop}[{\cite[Proposition 5.3]{BMM}}] %\label{prop:homological_block_simple_expression}
	\begin{gather*}
	\widehat{Z}_\Gamma (q) = \frac{1}{2} q^{(-18 - w_1 - \cdots - w_6 - 1/w_3 - 1/w_4 - 1/w_5 - 1/w_6)/4}
	\left( F_{Q_{+}, \veps}(\tau) - F_{Q_{-}, \veps}(\tau) \right).
	\end{gather*}
\end{prop}

